\documentclass[10pt]{article}
\usepackage[T1]{fontenc}
\usepackage[utf8]{inputenc}
\usepackage[french]{babel}
\usepackage{amsmath,amssymb,mathrsfs,mathpazo}
\usepackage{geometry}
\geometry{a4paper,margin=2cm}
\usepackage{graphicx,array,tabularx,tabularray,colortbl,multirow,multicol,booktabs,enumitem,pifont,ulem,fancybox,diagbox,cancel}
\usepackage{pstricks,pstricks-add,pst-plot,pst-tree,pst-node,pst-text,pst-func,pst-math,pst-bspline,pst-3dplot,pst-eucl,pst-solides3d}
\let\pscalculate\undefined
\usepackage{pst-bezier}
\usepackage{tikz,pgfplots,tkz-tab}
\usetikzlibrary{arrows,arrows.meta,patterns,calc,positioning,trees,decorations.markings,shapes.geometric,shapes.misc,shapes.arrows,intersections,decorations.pathmorphing,decorations.pathreplacing,backgrounds,shadows,babel}
\pgfplotsset{compat=1.18}
\usepgfplotslibrary{fillbetween}
\usepackage[np]{numprint}
\usepackage{listings,makecell,pas-tableur,twoopt}
\let\barre\undefined
\usepackage[pstricks]{tablvar}
\definecolor{darkgreen}{RGB}{0,100,0}
\providecommand{\red}{\color{red}}
\providecommand{\blue}{\color{blue}}
\providecommand{\green}{\color{green}}
\providecommand{\black}{\color{black}}
\providecommand{\white}{\color{white}}
\providecommand{\gray}{\color{gray}}
\providecommand{\cyan}{\color{cyan}}
\providecommand{\magenta}{\color{magenta}}
\providecommand{\yellow}{\color{yellow}}
\providecommand{\violet}{\color{violet}}
\providecommand{\orange}{\color{orange}}
\providecommand{\euro}{\texteuro}
\providecommand{\EUR}{\texteuro}
\providecommand{\ieme}{\textsuperscript{e}}
\providecommand{\degres}{\ensuremath{^\circ}}
\providecommand{\degree}{\ensuremath{^\circ}}
\providecommand{\vv}[1]{\overrightarrow{#1}}
\providecommand{\Coord}[1]{#1}
\providecommand{\decosix}{}
\providecommand{\starredbullet}{\ensuremath{\bullet}}
\pagestyle{empty}
\setlength{\parindent}{0pt}
\usepackage{adjustbox}\DeclareUnicodeCharacter{2236}{:}\DeclareUnicodeCharacter{2190}{\ensuremath{\leftarrow}}
\providecommand{\R}{\mathbb{R}}
\providecommand{\N}{\mathbb{N}}
\providecommand{\Z}{\mathbb{Z}}
\providecommand{\Q}{\mathbb{Q}}
\providecommand{\D}{\mathbb{D}}
\providecommand{\C}{\mathbb{C}}
\providecommand{\e}{{\mathrm{e}}}
\providecommand{\vect}[1]{\overrightarrow{#1}}
\providecommand{\vectt}[1]{\overrightarrow{\mathrm{#1}}}
\providecommand{\vvt}[1]{\overrightarrow{\mathrm{#1}}}
\providecommand{\np}[2][]{{#2~#1}}
\providecommand{\vv}[1]{{\overrightarrow{#1}}}
\renewcommand{\d}{{\mathrm{d}}}
\providecommand{\degres}{{^\circ}}
\providecommand{\degre}{{^\circ}}
\providecommand{\blue}{\color{blue}}
\providecommand{\red}{\color{red}}
\providecommand{\enskip}{\quad}
\providecommand{\llbracket}{\lbrack\!\lbrack}
\providecommand{\rrbracket}{\rbrack\!\rbrack}
\providecommand{\euro}{{\text{€}}}
\providecommand{\no}{{\text{n°}}}
\providecommand{\interval}[2]{{[#1;#2]}}
\providecommand{\intervFO}[2]{{[#1;#2[}}
\providecommand{\dotfill}{\ldots}

\providecommand{\ds}{\displaystyle}
\providecommand{\textding}[1]{#1}
\providecommand{\Oij}{\(\mathrm{O};\vec{\imath},\vec{\jmath}\)}
\providecommand{\Oijk}{\(\mathrm{O};\vec{\imath},\vec{\jmath},\vec{k}\)}
\providecommand{\Ouv}{\(\mathrm{O};\vec{u},\vec{v}\)}
\providecommand{\intervFF}[2]{\ensuremath{[#1;#2]}}
\renewcommand{\euro}{\ifmmode\text{\texteuro}\else\texteuro\fi}
\DeclareUnicodeCharacter{00B0}{\ifmmode^{\circ}\else\textdegree\fi}
\DeclareUnicodeCharacter{20AC}{\ifmmode\text{\texteuro}\else\texteuro\fi}
% Énoncés : Éducation nationale. Transcription et corrigés : APMEP, Denis Vergès et contributeurs. Découpage et classement : CoopMaths, Sébastien Lozano, Éric Elter et contributeurs. Adaptation de la présentation : Maths Explorer.
% CC BY-SA : https://coopmaths.fr/www/about/
% https://coopmaths.fr/alea/static/bac/2026/tex/bac_2026_06_sujet1_antilles_4.tex
\begin{document}
\noindent\textit{APMEP / CoopMaths. CC BY-SA. Présentation adaptée par Maths Explorer.}\par\medskip
\subsection*{Partie A : lecture graphique}

On considère une fonction $f$ définie et deux fois dérivable sur l'intervalle $\left]\dfrac{1}{2}~;~+\infty\right[$.
On note $f'$ sa fonction dérivée et $f''$ sa fonction dérivée seconde.

On a tracé, dans un repère orthogonal, la courbe de la fonction $f$, notée $\mathcal{C}$, ainsi que la tangente $T$ à $\mathcal{C}$ au point A(1~;~0). On précise que $T$ passe par le point B$(0~;~-2)$.

On répondra aux questions suivantes en les justifiant à l'aide du graphique.

\begin{center}




























\psset{yunit=0.8cm,xunit=2.25cm,arrowsize=2pt 3}
\begin{pspicture*}(-0.25,-3)(4.25,6.5)
\psgrid[gridlabels=0pt,subgriddiv=1,gridwidth=0.2pt]
\psaxes[linewidth=1.25pt, Dx=10,Dy=10]{->}(0,0)(-0.25,-3)(4.25,6.5)
\psplot[plotpoints=2000,linewidth=1.25pt,linecolor=blue]{0.501}{4}{2 x mul 1 sub ln x mul}
\psplotTangent[linestyle=dashed]{1}{5}{2 x mul 1 sub ln}
\uput[ul](1,0){A}\uput[d](1,0){1}\uput[l](0,1){1}
\uput[ul](3.2,5.5){\blue $\mathcal{C}$}\uput[r](3.8,5.2){$T$}\uput[ul](0,-2){B}\uput[dl](0,0){O}
\end{pspicture*}

\end{center}

\begin{enumerate}
\item Déterminer $f'(1)$.
\item La fonction $f$ est-elle concave sur $\left]\dfrac{1}{2}~;+\infty\right[$~?
\item Déterminer $f''(1)$.
\end{enumerate}


\subsection*{Partie B : étude de fonction}

On admet que la fonction représentée graphiquement dans la partie A est définie sur l'intervalle $\left]\dfrac{1}{2}~;~+\infty\right[$ par~:

\[f(x) = x\ln(2x-1)\]
où $\ln$ désigne la fonction logarithme népérien.

On rappelle que~:
\begin{itemize}[noitemsep]
\item la fonction $f$ est deux fois dérivable sur l'intervalle $\left]\dfrac{1}{2}~;+\infty\right[$, que $f'$ désigne sa fonction dérivée et $f''$ sa fonction dérivée seconde~;
\item $\mathcal{C}$ désigne la courbe représentative de la fonction $f$ sur l'intervalle $\left]\dfrac{1}{2}~;~+\infty\right[$.
\end{itemize}

\begin{enumerate}
\item Déterminer, en justifiant, la limite de la fonction $f$ en $\dfrac{1}{2}$. En déduire une interprétation graphique.
\item Résoudre l'équation $f(x) = 0$ sur l'intervalle $\left]\dfrac{1}{2}~;~+\infty\right[$.
\item On admet que pour tout réel $x$ appartenant à l'intervalle $\left]\dfrac{1}{2}~;+\infty\right[$~:

\[f'(x) = \ln(2x-1) + \dfrac{2x}{2x-1}.\]

	\begin{enumerate}[label=\textbf{\alph*.}]
		\item Montrer que pour tout réel $x$ appartenant à l'intervalle $\left]\dfrac{1}{2}~;+\infty\right[$~:

\[f''(x) = \dfrac{4(x-1)}{(2x-1)^2}.\]

		\item Montrer que $\mathcal{C}$ admet un unique point d'inflexion et que ce point est l'intersection de $\mathcal{C}$ avec l'axe des abscisses.
	\end{enumerate}
\item
	\begin{enumerate}[label=\textbf{\alph*.}]
		\item Montrer que la fonction $f$ est positive sur l'intervalle $[1~;~2]$.
		\item On admet que pour tout $x$ appartenant à l'intervalle $\left]\dfrac{1}{2}~;~+\infty\right[$, on a~:
		
\[\dfrac{x^2}{2x-1} = \dfrac{1}{2}x + \dfrac{1}{4} + \dfrac{1}{8}\times\dfrac{2}{2x-1}.\]

Montrer que~:

\[\displaystyle\int_{1}^{2} \dfrac{x^2}{2x-1}\, \text{d}x = 1 + \dfrac{1}{8}\ln (3).\]

\item On note $\mathcal{A}$ l'aire du domaine délimité par la courbe $\mathcal{C}$, l'axe des abscisses et les droites d'équation $x=1$ et $x=2$.

En vous aidant d'une intégration par parties, déterminer la valeur exacte de l'aire $\mathcal{A}$ en unité d'aire.
	\end{enumerate}
\end{enumerate}


\subsection*{Partie C : généralisation}

On désigne par $a$ un réel strictement positif et par $f_a$ la fonction définie sur l'intervalle $\left]\dfrac{1}{a}~;~+\infty\right[$ par~:

\[f_a(x) = x\ln(ax - 1).\]

On admet que $f_a$ est deux fois dérivable sur l'intervalle $\left]\dfrac{1}{a}~;~+\infty\right[$. On note $f_a'$ sa fonction dérivée, $f_a''$ sa fonction dérivée seconde, et $\mathcal{C}_a$ sa représentation graphique.

On admet que pour tout $x$ appartenant à l'intervalle $\left]\dfrac{1}{a}~;~+\infty\right[$:

\[f_a''(x) = \dfrac{a(ax-2)}{(ax-1)^2}.\]

\begin{enumerate}
\item Montrer que pour tout réel $a$ strictement positif, la courbe $\mathcal{C}_a$ admet un unique point d'inflexion A$_a$.
\item Montrer que les points A$_a$ sont alignés, $a$ appartenant à l'intervalle $]0~;~+\infty[$.
\end{enumerate}

\end{document}
